\documentclass{article}
\usepackage[utf8]{inputenc}
\begin{document}

\section{Codec Evaluation Methodology}

We follow the codec benchmarking methodology of Naderi
et~al.~\cite{naderi_vcd_2024}, encoding each source clip at five
quantization-parameter levels (QP\,=\,20--44 for H.264/H.265 and QP\,=\,22--42 for H.266).
Bjontegaard delta rate (BD-rate)
is computed per clip relative to the H.264 baseline using PSNR and VMAF
as quality metrics.

Recent work by Zhang et~al.~\cite{zhang_neural_2023} proposed a neural
codec that jointly optimizes rate, distortion, and perceptual quality
through a three-objective loss balancing strategy.

\subsubsection{Stratification strategy.}
\label{sec:stratification}

The subset selection employs a stratified sampling algorithm following Naderi
et~al.~\cite{naderi_vcd_2024} to ensure that each group covers the full
range of quality levels, spatial--temporal complexity, and distortion types.
Clips are assigned to one of 12 strata formed by the Cartesian product of
three MOS bins (Low: $[1.0, 2.8)$, Medium: $[2.8, 4.0)$, High:
$[4.0, 5.0]$) and four SI$\times$TI quadrants (split at the population
medians of SI and TI, computed per ITU-T
P.910~\cite{itu-t_recommendation_p910_subjective_2023}).

\bibliographystyle{plain}
\bibliography{sample_refs}

\end{document}
